\section{Distributions}

\begin{definition}
    Let $X$ be a random variable with values in $\{0,1\}$ and density
    $$P[X=1] = p \quad \text{and} \quad P[X=0] = 1-p$$
    Then $X$ follows a \textbf{Bernoulli distribution} with parameter $p$, denoted as $X \sim \text{Bernoulli}(p)$.
\end{definition}

\begin{example}
    Throwing a coin, indicatorvariable for heads.
\end{example}

\begin{remark}
    Observe that for a bernoulli-distribution we have:
    $$\mathbb{E}[X] = p \quad \text{and} \quad Var[X] = p(1-p)$$
\end{remark}

\begin{definition}
    Let $X$ be a random variable with values in $\{0,1,\dots,n\}$ and density
    $$P[X=k] = \binom{n}{k} p^k (1-p)^{n-k}$$
    Then $X$ follows a \textbf{Binomial distribution} with parameters $n$ and $p$, denoted as $X \sim \text{Bin}(n,p)$.
\end{definition}

\begin{example}
    Number of heads after throwing a coin $n$ times.
\end{example}

\begin{remark}
    Observe that for a binomial-distribution we have:
    $$\mathbb{E}[X] = np \quad \text{and} \quad Var[X] = np(1-p)$$
\end{remark}

\begin{definition}
    Let $X$ be a random variable with density
    $$P[X=k] = \frac{e^{-\lambda} \lambda^k}{k!}$$
    Then $X$ follows a \textbf{Poisson distribution} with parameter $\lambda$, denoted as $X \sim \text{Po}(\lambda)$.
\end{definition}

\begin{remark}
    Observe that for a poisson-distribution we have:
    $$\mathbb{E}[X] = \lambda \quad \text{and} \quad Var[X] = \lambda$$
\end{remark}

\begin{remark}
    This is often used to model events that do not happen often. For example the number of emails you get in a day. ($\lambda$ would be the average num of emails you get a day).
\end{remark}

\begin{definition}
    Let $X$ be a random variable with density
    $$P[X=k] = p(1-p)^{k-1}$$
    Then $X$ follows a \textbf{Geometric distribution} with parameter $p$, denoted as $X \sim \text{Geo}(p)$.
\end{definition}

\begin{example}
    Assume we throw a coin with $P[\text{head}] = p$ until we get n times heads. Let $X$ be the number of throws until we get n times heads.
    $$P[X=k] = \binom{k-1}{n-1} p^n (1-p)^{k-n}$$
\end{example}

\begin{remark}
    Observe that for a geometric-distribution we have:
    $$\mathbb{E}[X] = \frac{1}{p} \quad \text{and} \quad Var[X] = \frac{1-p}{p^2}$$
\end{remark}

\begin{definition}
    For $n \geq 2$ a random variable $X$ with density
    $$P[X=k] = \binom{k-1}{n-1} p^n (1-p)^{k-n}$$
    is called negative binomial distributed with order $n$.
\end{definition}

\begin{example}
    There are $n$ diffrent types of collectible cards. In every round we uniformly get one of them. Let $X$ be the number of rounds until we have all $n$ types of cards. What is $\mathbb{E}[X]$?
\end{example}